% assets/w09-role-split.svg — 9주차의 역할 분담 : 유도는 둘, 제어는 하나
%
%   bash _tools/tikz2svg.sh assets/src/w09-role-split.tex assets/w09-role-split.svg
%
% 이 그림이 말하는 것 하나
%   제어기는 5주차 것 그대로다. 바뀐 자리는 **제어기에 들어가는 명령을 누가 내는가**
%   하나뿐이고, 그 자리가 AvoidSeam 이다. 토픽 다섯이 두 세계를 잇는다.
%
% 유의 사항 둘 — 둘 다 2026-10-27 작업에서 실제로 발생한 오류다
%   (1) 줄바꿈은 두 겹 역슬래시만 쓴다. 길이를 달면 kotex/standalone 조합에서 수식
%       시작으로 읽혀 "Command \scriptsize invalid in math mode" 가 난다
%   (2) 이 파일은 **Write 도구로** 쓴다. 셸 히어독을 거치면 두 겹 역슬래시가 한 겹으로
%       줄어 줄바꿈이 통째로 사라진다 — 오류 없이 한 줄로 붙어 나오므로 더 나쁘다

\documentclass[border=6pt]{standalone}
% capstone-style.tex — 이 볼트의 TikZ 그림이 공유하는 스타일
%
% 쓰는 법:  % capstone-style.tex — 이 볼트의 TikZ 그림이 공유하는 스타일
%
% 쓰는 법:  % capstone-style.tex — 이 볼트의 TikZ 그림이 공유하는 스타일
%
% 쓰는 법:  \input{capstone-style}   (assets/src/ 안에서)
% 빌드:     bash _tools/tikz2svg.sh assets/src/w02-node-topic.tex assets/w02-node-topic.svg
%
% 왜 TikZ 인가
%   손으로 SVG 좌표를 쓰면 화살촉·선 굵기·글자 크기가 그림마다 미묘하게 달라진다.
%   그것이 "자료가 어설퍼 보인다"의 정체다. 스타일을 한 번만 정하고 상대좌표로
%   배치하면 좌표를 고쳐도 그림이 무너지지 않는다.
%   수식은 본문과 같은 TeX 글꼴로 나온다.
%
% 한글
%   ko.TeX(DVI 경로)를 쓴다. XeLaTeX 은 TikZ 그래픽을 PDF special 로 내보내는데
%   dvisvgm 이 그것을 받으려면 Ghostscript < 10.01 이 필요하다. 이 PC 에는
%   10.07 이 깔려 있어 그 경로가 막힌다. latex -> DVI -> dvisvgm 이 유일하게
%   동작하는 조합이다.

\usepackage[hangul]{kotex}
\usepackage{tikz}
\usepackage{amsmath}
\usepackage{xcolor}
\usetikzlibrary{arrows.meta, positioning, calc, fit, backgrounds, decorations.pathmorphing}

% ---- 색 --------------------------------------------------------------------
% 강조는 하나만. 나머지는 검정. 흑백 인쇄에서도 읽혀야 한다.
\definecolor{ink}{HTML}{1A1A1A}
\definecolor{soft}{HTML}{666666}
\definecolor{accent}{HTML}{7C3AED}
\definecolor{warn}{HTML}{B91C1C}

% 계층 색 — Simulink 모델의 블록 색과 같은 규칙을 쓴다
\definecolor{guid}{HTML}{1D4ED8}    % 유도
\definecolor{ctrl}{HTML}{EA580C}    % 제어
\definecolor{thr}{HTML}{D97706}     % 추진기
\definecolor{plant}{HTML}{16A34A}   % 운동모델
\definecolor{nav}{HTML}{7C3AED}     % 항법 · 신호처리
\definecolor{pathc}{HTML}{0D9488}   % 계획 경로 · 항적
\definecolor{errc}{HTML}{C2410C}    % 오차
\definecolor{flow}{HTML}{0369A1}    % 조류 · 바람

% ---- 선과 화살촉 -----------------------------------------------------------
\tikzset{
  sig/.style   = {ink, line width=0.5pt, -{Stealth[length=4pt, width=3pt]}},
  wire/.style  = {ink, line width=0.5pt},
  fb/.style    = {ink, line width=0.5pt, densely dashed,
                  -{Stealth[length=4pt, width=3pt]}},
  asig/.style  = {accent, line width=0.5pt, densely dashed,
                  -{Stealth[length=4pt, width=3pt]}},
  % 블록 — 흰 바탕, 직각, 검은 테두리. 색은 테두리에만 준다
  blk/.style   = {draw=ink, line width=0.55pt, fill=white,
                  minimum width=18mm, minimum height=8.5mm, inner sep=2pt, font=\small},
  gblk/.style  = {blk, draw=guid},
  cblk/.style  = {blk, draw=ctrl},
  tblk/.style  = {blk, draw=thr},
  pblk/.style  = {blk, draw=plant},
  nblk/.style  = {blk, draw=nav},
  % 합산점 · 분기점
  sum/.style   = {draw=ink, line width=0.55pt, fill=white, circle,
                  inner sep=0pt, minimum size=4.6mm, font=\scriptsize},
  dot/.style   = {ink, fill=ink, circle, inner sep=0pt, minimum size=1.5mm},
  % 글자
  sname/.style = {font=\small, inner sep=1.5pt},
  note/.style  = {font=\scriptsize, soft, inner sep=1.5pt},
  wnote/.style = {font=\scriptsize, warn, inner sep=1.5pt},
  anote/.style = {font=\scriptsize, accent, inner sep=1.5pt},
  axis/.style  = {ink, line width=0.5pt},
  tick/.style  = {font=\scriptsize, soft},
  % 묶음 상자 — 서브시스템 표시
  grp/.style   = {draw=soft, line width=0.4pt, densely dotted, rounded corners=1mm,
                  inner sep=3mm},
}
   (assets/src/ 안에서)
% 빌드:     bash _tools/tikz2svg.sh assets/src/w02-node-topic.tex assets/w02-node-topic.svg
%
% 왜 TikZ 인가
%   손으로 SVG 좌표를 쓰면 화살촉·선 굵기·글자 크기가 그림마다 미묘하게 달라진다.
%   그것이 "자료가 어설퍼 보인다"의 정체다. 스타일을 한 번만 정하고 상대좌표로
%   배치하면 좌표를 고쳐도 그림이 무너지지 않는다.
%   수식은 본문과 같은 TeX 글꼴로 나온다.
%
% 한글
%   ko.TeX(DVI 경로)를 쓴다. XeLaTeX 은 TikZ 그래픽을 PDF special 로 내보내는데
%   dvisvgm 이 그것을 받으려면 Ghostscript < 10.01 이 필요하다. 이 PC 에는
%   10.07 이 깔려 있어 그 경로가 막힌다. latex -> DVI -> dvisvgm 이 유일하게
%   동작하는 조합이다.

\usepackage[hangul]{kotex}
\usepackage{tikz}
\usepackage{amsmath}
\usepackage{xcolor}
\usetikzlibrary{arrows.meta, positioning, calc, fit, backgrounds, decorations.pathmorphing}

% ---- 색 --------------------------------------------------------------------
% 강조는 하나만. 나머지는 검정. 흑백 인쇄에서도 읽혀야 한다.
\definecolor{ink}{HTML}{1A1A1A}
\definecolor{soft}{HTML}{666666}
\definecolor{accent}{HTML}{7C3AED}
\definecolor{warn}{HTML}{B91C1C}

% 계층 색 — Simulink 모델의 블록 색과 같은 규칙을 쓴다
\definecolor{guid}{HTML}{1D4ED8}    % 유도
\definecolor{ctrl}{HTML}{EA580C}    % 제어
\definecolor{thr}{HTML}{D97706}     % 추진기
\definecolor{plant}{HTML}{16A34A}   % 운동모델
\definecolor{nav}{HTML}{7C3AED}     % 항법 · 신호처리
\definecolor{pathc}{HTML}{0D9488}   % 계획 경로 · 항적
\definecolor{errc}{HTML}{C2410C}    % 오차
\definecolor{flow}{HTML}{0369A1}    % 조류 · 바람

% ---- 선과 화살촉 -----------------------------------------------------------
\tikzset{
  sig/.style   = {ink, line width=0.5pt, -{Stealth[length=4pt, width=3pt]}},
  wire/.style  = {ink, line width=0.5pt},
  fb/.style    = {ink, line width=0.5pt, densely dashed,
                  -{Stealth[length=4pt, width=3pt]}},
  asig/.style  = {accent, line width=0.5pt, densely dashed,
                  -{Stealth[length=4pt, width=3pt]}},
  % 블록 — 흰 바탕, 직각, 검은 테두리. 색은 테두리에만 준다
  blk/.style   = {draw=ink, line width=0.55pt, fill=white,
                  minimum width=18mm, minimum height=8.5mm, inner sep=2pt, font=\small},
  gblk/.style  = {blk, draw=guid},
  cblk/.style  = {blk, draw=ctrl},
  tblk/.style  = {blk, draw=thr},
  pblk/.style  = {blk, draw=plant},
  nblk/.style  = {blk, draw=nav},
  % 합산점 · 분기점
  sum/.style   = {draw=ink, line width=0.55pt, fill=white, circle,
                  inner sep=0pt, minimum size=4.6mm, font=\scriptsize},
  dot/.style   = {ink, fill=ink, circle, inner sep=0pt, minimum size=1.5mm},
  % 글자
  sname/.style = {font=\small, inner sep=1.5pt},
  note/.style  = {font=\scriptsize, soft, inner sep=1.5pt},
  wnote/.style = {font=\scriptsize, warn, inner sep=1.5pt},
  anote/.style = {font=\scriptsize, accent, inner sep=1.5pt},
  axis/.style  = {ink, line width=0.5pt},
  tick/.style  = {font=\scriptsize, soft},
  % 묶음 상자 — 서브시스템 표시
  grp/.style   = {draw=soft, line width=0.4pt, densely dotted, rounded corners=1mm,
                  inner sep=3mm},
}
   (assets/src/ 안에서)
% 빌드:     bash _tools/tikz2svg.sh assets/src/w02-node-topic.tex assets/w02-node-topic.svg
%
% 왜 TikZ 인가
%   손으로 SVG 좌표를 쓰면 화살촉·선 굵기·글자 크기가 그림마다 미묘하게 달라진다.
%   그것이 "자료가 어설퍼 보인다"의 정체다. 스타일을 한 번만 정하고 상대좌표로
%   배치하면 좌표를 고쳐도 그림이 무너지지 않는다.
%   수식은 본문과 같은 TeX 글꼴로 나온다.
%
% 한글
%   ko.TeX(DVI 경로)를 쓴다. XeLaTeX 은 TikZ 그래픽을 PDF special 로 내보내는데
%   dvisvgm 이 그것을 받으려면 Ghostscript < 10.01 이 필요하다. 이 PC 에는
%   10.07 이 깔려 있어 그 경로가 막힌다. latex -> DVI -> dvisvgm 이 유일하게
%   동작하는 조합이다.

\usepackage[hangul]{kotex}
\usepackage{tikz}
\usepackage{amsmath}
\usepackage{xcolor}
\usetikzlibrary{arrows.meta, positioning, calc, fit, backgrounds, decorations.pathmorphing}

% ---- 색 --------------------------------------------------------------------
% 강조는 하나만. 나머지는 검정. 흑백 인쇄에서도 읽혀야 한다.
\definecolor{ink}{HTML}{1A1A1A}
\definecolor{soft}{HTML}{666666}
\definecolor{accent}{HTML}{7C3AED}
\definecolor{warn}{HTML}{B91C1C}

% 계층 색 — Simulink 모델의 블록 색과 같은 규칙을 쓴다
\definecolor{guid}{HTML}{1D4ED8}    % 유도
\definecolor{ctrl}{HTML}{EA580C}    % 제어
\definecolor{thr}{HTML}{D97706}     % 추진기
\definecolor{plant}{HTML}{16A34A}   % 운동모델
\definecolor{nav}{HTML}{7C3AED}     % 항법 · 신호처리
\definecolor{pathc}{HTML}{0D9488}   % 계획 경로 · 항적
\definecolor{errc}{HTML}{C2410C}    % 오차
\definecolor{flow}{HTML}{0369A1}    % 조류 · 바람

% ---- 선과 화살촉 -----------------------------------------------------------
\tikzset{
  sig/.style   = {ink, line width=0.5pt, -{Stealth[length=4pt, width=3pt]}},
  wire/.style  = {ink, line width=0.5pt},
  fb/.style    = {ink, line width=0.5pt, densely dashed,
                  -{Stealth[length=4pt, width=3pt]}},
  asig/.style  = {accent, line width=0.5pt, densely dashed,
                  -{Stealth[length=4pt, width=3pt]}},
  % 블록 — 흰 바탕, 직각, 검은 테두리. 색은 테두리에만 준다
  blk/.style   = {draw=ink, line width=0.55pt, fill=white,
                  minimum width=18mm, minimum height=8.5mm, inner sep=2pt, font=\small},
  gblk/.style  = {blk, draw=guid},
  cblk/.style  = {blk, draw=ctrl},
  tblk/.style  = {blk, draw=thr},
  pblk/.style  = {blk, draw=plant},
  nblk/.style  = {blk, draw=nav},
  % 합산점 · 분기점
  sum/.style   = {draw=ink, line width=0.55pt, fill=white, circle,
                  inner sep=0pt, minimum size=4.6mm, font=\scriptsize},
  dot/.style   = {ink, fill=ink, circle, inner sep=0pt, minimum size=1.5mm},
  % 글자
  sname/.style = {font=\small, inner sep=1.5pt},
  note/.style  = {font=\scriptsize, soft, inner sep=1.5pt},
  wnote/.style = {font=\scriptsize, warn, inner sep=1.5pt},
  anote/.style = {font=\scriptsize, accent, inner sep=1.5pt},
  axis/.style  = {ink, line width=0.5pt},
  tick/.style  = {font=\scriptsize, soft},
  % 묶음 상자 — 서브시스템 표시
  grp/.style   = {draw=soft, line width=0.4pt, densely dotted, rounded corners=1mm,
                  inner sep=3mm},
}


\begin{document}
\begin{tikzpicture}[x=1mm, y=1mm]

% ===================== Simulink 사슬 (위) =====================
\node[gblk, text width=24mm, align=center] (los)  at ( 22,90)
     {LOS 유도 \\ {\scriptsize 5주차 그대로}};
\node[gblk, text width=26mm, align=center] (seam) at ( 58,90)
     {\textbf{AvoidSeam} \\ {\scriptsize 어느 명령을 쓸까}};
\node[cblk, text width=30mm, align=center] (ctrl) at ( 98,90)
     {헤딩 P-D · 속도 PI \\ {\scriptsize 한 줄도 안 바뀜}};
\node[tblk, text width=20mm, align=center] (alloc) at (134,90) {추력배분};
\node[nblk, text width=28mm, align=center] (goal)  at ( 22,68) {GoalPublisher};

\draw[sig] (los.east)  -- (seam.west);
\draw[sig] (seam.east) -- (ctrl.west);
\draw[sig] (ctrl.east) -- (alloc.west);
\node[note, anchor=south] at (40,91.5) {$\psi_{los}$};
\node[note, anchor=south] at (78,91.5) {$\psi_{ref}$};

\node[grp, fit=(los)(seam)(ctrl)(alloc)(goal)] (sim) {};
\node[note, anchor=north west] at (sim.north west) {\textbf{Simulink} — 제어는 여기 하나뿐};

% ===================== 파이썬 (아래 왼쪽) =====================
\node[nblk, text width=54mm, align=center] (avoid) at (38,48)
     {회피기 \texttt{collision\_avoidance\_enu} \\ {\scriptsize 호 다발 21개 중 하나를 고른다}};
\node[nblk, text width=30mm, align=center] (bridge) at (38,28) {\texttt{ca\_bridge}};

\node[grp, fit=(avoid)(bridge)] (py) {};
\node[note, anchor=south west] at (py.north west) {\textbf{ROS 2 · Python} — 유도 하나를 더 낸다};

% ===================== Gazebo (아래 오른쪽) =====================
\node[pblk, text width=30mm, align=center] (lidar) at (120,48)
     {3D LiDAR \\ {\scriptsize $z_b=1.90$ m · $16\times440$}};
\node[pblk, text width=30mm, align=center] (odom)  at (120,28) {위치 · 자세};
\node[tblk, text width=20mm, align=center] (thrb)  at (154,38) {추진기 셋};

\node[grp, fit=(lidar)(odom)(thrb)] (gz) {};
\node[note, anchor=south west] at (gz.north west) {\textbf{Gazebo · VRX} \texttt{stationkeeping\_task}};

% ===================== 센서 둘 =====================
\draw[sig] (lidar.west) -- (avoid.east);
\node[note, anchor=south] at (85,49.5) {\texttt{lidar\_wamv\_sensor/scan}};

\draw[sig] (odom.west) -- (bridge.east);
\node[note, anchor=south] at (81,29.5) {\texttt{ground\_truth\_odometry}};

\draw[sig] (bridge.north) -- (avoid.south);
\node[note, anchor=west] at (39.5,38) {\texttt{/wamv/gps\_pose}};

% ===================== 이음새로 올라가는 토픽 둘 =====================
\draw[sig] (58,52.5) -- (seam.south);
\node[note, anchor=west, text width=40mm] at (59.5,67)
     {\texttt{/vfh/command} \\ angular.z $=$ 절대 방위각 (ENU, rad) \\ ~ \\
      \texttt{/vfh/oa\_flag} \\ 피하는 중이면 true};

% ===================== 목표점은 거꾸로 내려온다 =====================
\draw[fb] (los.south) -- (goal.north);
\node[note, anchor=west] at (23.5,79) {\texttt{wp\_idx}};
\draw[sig] (goal.south) -- (22,52.5);
% 이 꼬리표는 y=59 에 두면 아래 파이썬 묶음의 제목줄과 겹친다 (2026-10-27 에 겹쳤다)
\node[note, anchor=east] at (20.5,62) {\texttt{/usv/state/goal\_enu}};

% ===================== 위치는 양쪽이 같은 토픽에서 받는다 =====================
\draw[fb] (odom.east) -- (170,28) -- (170,105) -- (98,105) -- (ctrl.north);
\node[note, anchor=south east] at (168,106) {같은 토픽 — 자를 하나로 쓴다};

% ===================== 추력 =====================
\draw[sig] (134,82.5) -- (134,58) -- (154,58) -- (thrb.north);
% 오른쪽으로 뻗게 두면 x=170 되먹임선을 뚫고 그림 밖으로 나간다. 가로선 위에 얹는다
\node[note, anchor=south] at (144,58.8) {\texttt{/wamv/thrusters/*/thrust}};

% ===================== 그림은 파이썬이 그린다 =====================
\node[nblk, text width=42mm, align=center] (viz) at (84,6)
     {\texttt{ca\_viz} — 극좌표 실시간 그림};

\draw[fb] (14,43.5) -- (14,16) -- (84,16) -- (viz.north);
\node[note, anchor=west] at (38,17.5) {스캔 · 위치 · 회피 명령};

\draw[fb] (sim.west) -- (-4,80) -- (-4,6) -- (viz.west);
% 한 줄로 두면 ca_viz 상자를 뚫고 지나간다. 두 줄로 접어 상자 아래에 둔다
\node[note, anchor=north west, text width=60mm] at (-2,1)
     {\texttt{/sim/psi\_los} · \texttt{/sim/psi\_cmd} \\ 모델이 자기만 아는 값 둘을 낸다};

% ===================== 주의 =====================
\node[wnote, anchor=north west, text width=150mm] at (-4,-12)
     {\textbf{이 그림에서 5주차와 달라진 자리는 \texttt{AvoidSeam} 한 칸뿐이다.}
      헤딩 P-D · 속도 PI · 추력배분 · 운동모델은 블록 하나도 바뀌지 않았다.
      유도가 둘(LOS 와 회피)로 늘고, 그 둘 중 하나를 고르는 상자가 생긴 것이 전부다};

\end{tikzpicture}
\end{document}
